\documentclass[14pt, a4paper]{extarticle}
\usepackage[T2A]{fontenc}
\usepackage[utf8]{inputenc}
\usepackage[english, russian]{babel}
\usepackage[left=30mm, right=10mm, top=20mm, bottom=25mm]{geometry}

\usepackage{misccorr}
\usepackage{indentfirst}

% Заголовки section по центру
\makeatletter
\renewcommand\section{\@startsection {section}{1}{\z@}%
                                   {-3.5ex \@plus -1ex \@minus -.2ex}%
                                   {2.3ex \@plus.2ex}%
                                   {\centering\normalfont\Large\bfseries}}
\makeatother

% Точки в списке литературы
\makeatletter
\renewcommand\@biblabel[1]{#1.}
\makeatother

\begin{document}
\thispagestyle{empty}
{\scriptsize
\begin{center}
ФЕДЕРАЛЬНОЕ ГОСУДАРСТВЕННОЕ ОБРАЗОВАТЕЛЬНОЕ БЮДЖЕТНОЕ УЧРЕЖДЕНИЕ\\ 
ВЫСШЕГО ОБРАЗОВАНИЯ    
\end{center}
}
{\bf
\begin{center}
ФИНАНСОВЫЙ УНИВЕРСИТЕТ\\ 
ПРИ ПРАВИТЕЛЬСТВЕ РОССИЙСКОЙ ФЕДЕРАЦИИ\\
Липецкий филиал
\end{center}
}
\hrule
\bigskip

{\small 
\begin{center}
\textbf{
Кафедра <<Учет и информационные технологии в бизнесе>>}
\end{center}
}
\bigskip

Направление: Прикладная математика и информатика                                                                                                                                                                                                                                                                                                     

\medskip
Профиль: Анализ данных

\vskip 2cm

\begin{center}
\textbf{КУРСОВОЙ ПРОЕКТ}
\end{center}

по дисциплине: <<Машинное обучение>>

\bigskip
\begin{center}
ТЕМА: Анализ поведенческих закономерностей пользователей\\ социальных сетей
\end{center}

\vskip 3.0cm

\hskip 0.28\textwidth Преподаватель: к.ф.-м.н., доцент Калитвин В. А.

\hskip 0.28\textwidth Студент: Шуркалова А. Н.

\hskip 0.28\textwidth Личное дело \textnumero 10016240077

\hskip 0.28\textwidth Курс второй

\vskip 7.5cm

\begin{center}
Липецк 2025
\end{center}

\tableofcontents
\newpage

% ========== ВВЕДЕНИЕ ==========
\section*{ВВЕДЕНИЕ}
\addcontentsline{toc}{section}{ВВЕДЕНИЕ}

В настоящее время социальные сети стали неотъемлемой частью нашей жизни. Ежедневно мы оставляем там большие объёмы своих данных, 
так называемый цифровой след. Эта информация показывает поведенческие паттерны человека, то есть устойчивые модели поведения, 
поддающиеся выявлению и анализу. Знание этих закономерностей важно как для самих платформ, чтобы улучшить системы рекомендаций 
и удержать внимание пользователя. Также эта информация нужна и бизнесу, чтобы адаптировать рекламу под интересы читателя.

По некоторым исследованиям выявлено, что использование социальных сетей носит привычный характер, то есть поведение пользователя
можно предугадать. Как раз для выявления таких паттернов используются методы машинного обучения, в частности алгоритмы классификации
и кластеризации. Классификация --- это метод, который относит объект к одному из заранее известных классов (например, «активный пользователь» или «пассивный»). Кластеризация же позволяет группировать объекты без обучающих меток, определяя скрытые закономерности в данных.
 Также в российской науке активно развиваются методы анализа цифрового следа пользователей на основе данных социальных сетей.
Но вопрос о том, какие алгоритмы и метрики расстояния покажут наиболее точные результаты при анализе именно поведенческих паттернов, остаётся открытым.

В данной работе объектом исследования являются социальные сети как среда взаимодействия пользователей и генерация поведенческих данных.
Предметом исследования служат поведенческие закономерности пользователей, выраженные в числовых признаках, таких как частота
взаимодействия пользователей с контентом, время активности или другие действия, отражающие вовлечённость.

Целью курсовой работы является анализ поведенческих закономерностей пользователей социальных сетей с использованием методов машинного
обучения и выявление наиболее подходящего подхода к кластеризации и классификации пользователей на основе их паттернов.

Чтобы достигнуть поставленной цели, понадобится решить следующие задачи:

1. Изучить теоретические основы анализа социальных сетей и методов машинного обучения, а именно метрики расстояния, такие как евклидово,
манхэттенское, Чебышёва, и алгоритмы KNN (K‑ближайших соседей), K‑Means.

2. Рассмотреть, какие есть подходы к анализу поведенческих паттернов пользователей в социальных сетях, и выделить ключевые поведенческие
признаки.

3. Найти и подготовить датасет, который будет содержать в себе числовые данные о поведении пользователей, пригодные для применения
методов KNN и K‑Means.

4. Реализовать алгоритмы KNN (для классификации) и K‑Means (для кластеризации) с использованием языка Python и различных библиотек.

5. Провести эксперименты вычислений: подобрать оптимальные параметры моделей (значение K для KNN, число кластеров для K‑Means),
оценить качество классификации (точность, полнота, F1‑мера) и кластеризации (инерция, силуэт).

6. Истолковать результат, который получится, и сформулировать рекомендации по персонализации контента для различных групп пользователей.

После выполнения работы планируется получить модель, применимую на практике. С её помощью можно будет делить людей на группы по их 
поведению в соцсетях. Подразумевается, что метод K‑Means будет делить людей на несколько групп, с помощью этих групп станет легко 
выявить активных или пассивных пользователей. А алгоритм KNN, как предполагается, будет правильно определять принадлежность пользователя,
то есть достигнет точности классификации как минимум 80\% на тестовой выборке.

Курсовая работа состоит из введения, двух глав, заключения и списка литературы. В первой главе рассматриваются теоретические основы 
анализа социальных сетей и алгоритмов машинного обучения. Во второй главе приводится описание датасета, предобработка данных, реализация
моделей, анализ результатов и рекомендации.

% ========== ГЛАВА 1 ==========
\section{Общая характеристика алгоритмов и методов машинного обучения}
\subsection{Описание задачи и этапов её решения}
% ===== НОВАЯ ПОДГЛАВА (заглушка) =====
\textit{Текст будет написан после получения модели.}
\subsection{Анализ поведения пользователей в социальных сетях}

\textbf{Данные в социальных сетях}
\medskip

Данные соцсетей делятся на различный контент и связи между пользователями.
Контент максимально разнообразный, и для каждой социальной сети есть свои предпочтения:

\begin{itemize}
    \item с текстовым контентом (Twitter, Telegram);
    \item с преобладанием изображений (Instagram);
    \item с преобладанием видео (YouTube);
    \item смешанного типа (ВКонтакте, Facebook).
\end{itemize}

Зачастую данные соцсетей неструктурированные, и на это стоит обратить внимание, потому что потребуется
применение специальных методов обработки, таких как нормализация, стандартизация и выделение числовых признаков, которые в дальнейшем используются в машинном обучении в таких алгоритмах, как KNN и K‑Means.

\textbf{Поведенческие паттерны пользователей}
\medskip

По мнению авторов, анализ социальных сетей изучает поведение отдельных пользователей на микроуровне, общую
структуру связей на макроуровне и их взаимодействие~\cite{savin2020}.
В исследованиях активности пользователей обычно выделяют три типа:

\begin{itemize}
    \item Активные пользователи --- те, кто наиболее часто публикует различный контент, ставит лайки и взаимодействует с другими.
    \item Средние пользователи --- самая распространённая группа, иногда взаимодействуют с другими, реже что-то публикуют, ставят лайки.
    \item Пассивные пользователи --- зарегистрированы на площадке, но почти не проявляют активности.
\end{itemize}

Исследование пользователях в Индонезии выделило два кластера:

\begin{itemize}
    \item Кластер 0 (активные): пользуются Instagram, TikTok, YouTube по несколько часов в день, очень высокий уровень взаимодействия, интересуются развлекательным и мотивационным контентом.
    \item Кластер 1 (пассивные): в основном Instagram и TikTok, 3–4 часа в день, низкий уровень вовлечённости, предпочитают мотивационный контент и саморазвитие.
\end{itemize}

Это подтверждает, что K-Means эффективно и хорошо сегментирует пользователей по поведению.

\textbf{Поведенческие признаки, используемые для анализа}
\medskip

Чтобы оценить поведение пользователей, используют целевые действия, такие как лайки, репосты, комментарии,
переходы по ссылкам. Также учитывается время, проведённое в соцсети.
Кроме того, в датасетах для анализа часто применяют демографические признаки (возраст, пол, образование,
локацию), метрики вовлечённости (количество подписчиков, постов, время на платформах), а также активность
в будние и выходные дни.

Для определения активности пользователей исследователи использовали топологические признаки эго-сетей (структуры личных связей пользователя) и сделали вывод, что активные пользователи показывают более сложную динамику и разнообразные паттерны, чем пассивные.

\textbf{Задачи анализа поведения пользователей}
\medskip

На основе обзорной статьи по ML для соцсетей авторы отмечают, что соцсети являются одними из самых богатых источников своевременных и полных онлайн-данных. Методы машинного обучения позволяют:

\begin{itemize}
    \item Сегментировать аудиторию --- выделять группы пользователей со схожими паттернами поведения;
    \item Классифицировать пользователей --- определять, относится ли пользователь к активным или пассивным;
    \item Прогнозировать действия --- предсказывать, будет ли пользователь взаимодействовать с контентом;
    \item Персонализировать контент --- адаптировать рекомендации под конкретные группы пользователей.
\end{itemize}

Из исследования про Telegram авторы выделили 4 типа пользователей в дискуссиях о COVID-19 и изменении климата: Strategic Disruptor, Empirical Enthusiast, Inquisitive Moderate, Critical Examiner. Классификация достигла точности 0.95–0.99, то есть 95–99\%.

\textbf{Роль машинного обучения в анализе поведения}
\medskip

\begin{itemize}
    \item Кластеризация (K-Means) применяется, когда неизвестно, какие группы существуют. Алгоритм способен самостоятельно находить скрытые закономерности и объединять пользователей по общим интересам.
    \item Классификация (KNN) используется, когда есть размеченные примеры (например, «активный/пассивный»).
    \item Ключевая идея: даже относительно простые алгоритмы (как KNN) могут давать высокую точность при правильном выборе признаков и метрик расстояния.
\end{itemize}

Как отмечают авторы статьи об активности пользователей, несмотря на разнообразие поведения пользователей, метод может точно улавливать топологические особенности сети и эффективно классифицировать пользователей по уровню активности. Методы, основанные на машинном обучении, позволяют системам самостоятельно извлекать новые знания из обилия данных. За последние десятилетия многие статьи исследовали социальные сети через призму машинного обучения.

\subsection{Метрики расстояния}

Чтобы алгоритм мог определить, насколько два объекта похожи, необходимо уметь измерять расстояние между
ними. Чем меньше расстояние между двумя точками в пространстве признаков, тем более похожи эти объекты.

Существует большое разнообразие метрик расстояния, одной из самых распространённых является евклидово 
расстояние.

\textbf{Евклидово расстояние}
\medskip

Это расстояние по прямой между двумя точками. Для двух точек 
\( A(x_1, x_2, \dots, x_n) \) и \( B(y_1, y_2, \dots, y_n) \) 
оно вычисляется по формуле:

\[
d = \sqrt{(x_1 - y_1)^2 + \dots + (x_n - y_n)^2}
\]

Эта метрика часто используется для сравнения пользователей в многомерном пространстве признаков.
Например, если у нас есть два пользователя с известными параметрами «время в сети» и «количество лайков»,
евклидово расстояние покажет, насколько они похожи. Чем меньше расстояние, тем более схоже их поведение.

\textbf{Манхэттенское расстояние}
\medskip

Манхэттенское расстояние, также называемое расстоянием такси или расстоянием городского квартала, 
— это расстояние между двумя точками, которое равно сумме модулей разностей их координат.
В отличие от евклидова расстояния, манхэттенское расстояние измеряет дистанцию не по кратчайшей прямой,
а по блокам.

Формула манхэттенского расстояния:

\[
d = |x_1 - y_1| + |x_2 - y_2| + \dots + |x_n - y_n|
\]

В контексте анализа поведения пользователей манхэттенское расстояние может быть полезно, если признаки 
не связаны между собой и важны индивидуальные отклонения по каждой шкале. Например, при сравнении 
пользователей по таким показателям, как частота комментариев, количество репостов и время просмотра, 
эта метрика учитывает разницу по каждому признаку отдельно, не сглаживая её, в отличие от евклидова 
расстояния.

\textbf{Расстояние Чебышёва}
\medskip

Расстояние Чебышёва — это метрика на векторном пространстве, которая определяется как максимум модуля
разности соответствующих компонент двух векторов.

Формула расстояния Чебышёва:

\[
d = \max(|x_1 - y_1|, |x_2 - y_2|, \dots, |x_n - y_n|)
\]

Эта метрика ориентируется на самый «сильный» признак, по которому пользователи различаются больше всего. 
В задачах анализа поведения это может быть полезно, если один из поведенческих показателей (например, время в сети) 
значительно важнее остальных, и именно он определяет схожесть или различие пользователей.

\subsection{Алгоритм K-ближайших соседей}

Алгоритм KNN — это надежный и интуитивно понятный способ решения задач классификации и регрессии. Основной принцип алгоритма основан на нахождении K ближайших соседей на основе похожести их признаков.

Механизм K-ближайших соседей рассчитан на выявление расстояния между точками. При вводе новой точки запроса алгоритм выполняет поиск в сохраненном наборе данных, чтобы найти K обучающих выборок, наиболее близких к новому вводу.

\begin{enumerate}
    \item Измерение расстояния: наиболее распространенной мерой является евклидово расстояние, которое измеряет расстояние по прямой между точками в многомерном пространстве.
    
    \item Выбор соседей: после вычислений расстояний алгоритм сортирует их и определяет K ближайших записей.
    
    \item Принятие решений:
    \begin{itemize}
        \item Для классификации работает правило «Голосование большинством». Алгоритм смотрит на K ближайших соседей, определяет, к какому классу принадлежит каждый из них, и тот класс, который встречается чаще, побеждает.
        \item Для регрессии действует правило «усреднение». Алгоритм смотрит на K ближайших соседей, определяет числовое значение у каждого из них, считает среднее арифметическое этих значений и принимает решение.
    \end{itemize}
\end{enumerate}

\textbf{Преимущества и недостатки KNN}
\medskip

Как и любой алгоритм, KNN имеет свои сильные и слабые стороны.

\textbf{Преимущества KNN}
\medskip

\begin{itemize}
    \item Простота реализации и понимания. Алгоритм легко объяснить даже человеку, не знакомому с машинным обучением.
    \item Отсутствие этапа обучения. KNN не требует времени на обучение модели, так как просто запоминает все данные. За это свойство его называют «ленивым» алгоритмом.
    \item Универсальность. KNN может решать как задачи классификации, так и задачи регрессии.
    \item Устойчивость к выбросам. При правильном выборе параметра K влияние случайных выбросов снижается.
\end{itemize}

\textbf{Недостатки KNN}
\medskip

\begin{itemize}
    \item Медленная работа на больших данных. При каждом предсказании алгоритм вычисляет расстояния до всех объектов обучающей выборки, что требует много времени при большом объёме данных.
    \item Чувствительность к выбору K и метрики расстояния. Неправильный выбор параметров может существенно ухудшить качество предсказаний.
    \item Чувствительность к масштабу признаков. Если признаки имеют разные диапазоны значений, метрики расстояния могут искажаться. Поэтому перед применением KNN часто требуется нормализация данных.
    \item Проблема «проклятия размерности». При большом количестве признаков расстояния между объектами становятся почти одинаковыми, что снижает качество работы алгоритма.
\end{itemize}

\subsection{Алгоритм K-Means}

В машинном обучении существует также популярный алгоритм кластеризации — K-Means. Он относится к обучению без учителя: ему не нужны размеченные примеры. K-Means сам разбивает объекты на K групп (кластеров) так, чтобы объекты внутри одного кластера были максимально похожи, а объекты из разных кластеров — максимально различны.

В отличие от KNN, который предсказывает класс на основе известных данных, K-Means ищет скрытые группы там, где правильные ответы неизвестны.

\textbf{Как работает алгоритм K-Means}
\medskip

Алгоритм работает итеративно и состоит из следующих шагов:

\begin{enumerate}
    \item Выбор числа кластеров K — задаётся заранее (например, мы хотим разделить пользователей на 3 типа: пассивные, средние, активные).
    \item Инициализация центров — случайным образом выбираются K точек, которые становятся центрами кластеров.
    \item Распределение объектов — каждый объект (пользователь) относится к тому кластеру, центр которого находится ближе всего (расстояние обычно считается по евклидовой метрике).
    \item Пересчёт центров — после того как все объекты распределены, центры кластеров пересчитываются как среднее арифметическое всех точек, входящих в кластер.
    \item Повторение — шаги 3 и 4 повторяются до тех пор, пока центры кластеров не перестанут меняться, либо пока не будет выполнено заданное количество итераций.
\end{enumerate}

\textbf{Пример из жизни}
\medskip

Допустим, есть данные о пользователях: сколько времени они проводят в соцсети и сколько лайков ставят за день. Алгоритм K-Means может сам выделить группы:

\begin{itemize}
    \item Кластер 0 (пассивные): мало времени, мало лайков;
    \item Кластер 1 (средние): среднее время, среднее количество лайков;
    \item Кластер 2 (активные): много времени, много лайков.
\end{itemize}

При этом алгоритму не нужно говорить заранее, кто к какому типу относится. Он сам находит эти группы на основе данных.

\textbf{Преимущества K-Means}
\medskip

\begin{itemize}
    \item Простота — алгоритм легко понять и реализовать.
    \item Быстрота — работает быстро даже на больших наборах данных.
    \item Масштабируемость — хорошо справляется с большим количеством объектов.
    \item Интерпретируемость — результаты кластеризации легко анализировать, так как каждый объект попадает ровно в один кластер.
\end{itemize}

\textbf{Недостатки K-Means}
\medskip

\begin{itemize}
    \item Нужно заранее знать K — число кластеров надо задавать вручную. Подбирается методом локтя или по силуэту (об этом в параграфе 1.5).
    \item Чувствительность к начальным центрам — если неудачно выбрать начальные центры, алгоритм может сойтись к плохому результату. Проблема решается запуском алгоритма несколько раз с разными начальными условиями.
    \item Чувствительность к выбросам — один резко отличающийся пользователь может сильно сместить центр кластера.
    \item Только выпуклые кластеры — алгоритм плохо работает, если кластеры имеют сложную форму (например, в виде полумесяца или кольца).
\end{itemize}

\textbf{Применение K-Means к анализу поведения пользователей}
\medskip

В контексте данной работы K-Means используется для сегментации пользователей на основе их поведенческих признаков. Например, можно выделить группы пользователей по следующим параметрам:

\begin{itemize}
    \item время, проведённое в соцсети;
    \item количество лайков;
    \item количество комментариев;
    \item частота публикаций.
\end{itemize}

Полученные кластеры затем интерпретируются (например, «активные», «средние», «пассивные») и используются для персонализации контента или выработки рекомендаций.

\subsection{Методы оценки качества моделей}

Чтобы оценить, насколько качественно модель справилась с поставленной задачей, используют специальные метрики. Для задач классификации (KNN) и кластеризации (K-Means) применяются разные подходы~\cite{vorontsov}.

\textbf{Метрики для классификации (KNN)}
\medskip

При оценке качества классификации обычно используют следующие метрики~\cite{chubukova2024}:

\begin{itemize}
    \item Accuracy (точность) — это метрика, которая показывает долю правильно классифицированных объектов среди всех предсказанных. Другими словами, сколько объектов модель угадала верно из всех, что предсказала.
    
    \item Precision (точность предсказания положительного класса) — из всех объектов, которые модель отнесла к положительному классу (например, «активные пользователи»), сколько действительно являются таковыми.
    
    \item Recall (полнота) — из всех реально существующих объектов положительного класса, сколько модель смогла обнаружить.
    
    \item F1-мера — среднее гармоническое между Precision и Recall. Используется, когда нужно учесть обе метрики вместе.
    
    \item Матрица ошибок (confusion matrix) — таблица, показывающая количество правильных и ошибочных предсказаний. По диагонали располагаются верные ответы (TP и TN), а на побочной — ошибки (FP и FN).
\end{itemize}

TP (True Positive) — объекты правильно отнесённые к положительному классу. FN (False Negative) — объекты, которые модель ошибочно не отнесла к положительному классу. FP (False Positive) — объекты, которые модель ошибочно отнесла к положительному классу. TN (True Negative) — объекты правильно отнесённые к отрицательному классу~\cite{chubukova2024}.

Эти метрики дают полное представление о том, как модель справляется с задачей. Только точности (accuracy) бывает недостаточно, если классы в данных сильно несбалансированы.

\textbf{Метрики для кластеризации (K-Means)}
\medskip

Для оценки качества кластеризации, когда у нас нет правильных ответов, применяют внутренние метрики~\cite{chubukova2024}:

\begin{itemize}
    \item Инерция (inertia) — сумма квадратов расстояний от каждой точки до центра своего кластера. Чем меньше значение инерции, тем компактнее получились кластеры. Однако при увеличении числа кластеров инерция всегда уменьшается, поэтому её используют не как абсолютную меру, а для сравнения разных значений K.
    
    \item Силуэт (silhouette score) — показывает, насколько объект похож на свой кластер по сравнению с соседними. Значение силуэта лежит в диапазоне от -1 до 1. Чем ближе к 1, тем лучше: объект хорошо сгруппирован с «соседями» по кластеру и далёк от других кластеров. Значения около 0 означают, что кластеры пересекаются. Отрицательные значения говорят о том, что объект, вероятно, попал не в свой кластер.
\end{itemize}

\textbf{Как выбирать параметры моделей}
\medskip

Для KNN выбор оптимального значения K осуществляется с помощью кросс-валидации~\cite{vorontsov}. Набор данных разбивается на несколько частей (например, на 5 блоков), модель обучается на 4 из них и проверяется на 1. Так повторяется несколько раз, и в итоге выбирается тот K, при котором средняя точность максимальна.

Для K-Means число кластеров K подбирают, анализируя график инерции и значение силуэта:

\begin{itemize}
    \item Метод локтя — строится график зависимости инерции от числа кластеров. Оптимальное значение K находится в точке, где график резко замедляет своё падение, образуя «локоть».
    \item Максимум силуэта — перебираются разные значения K, и выбирается то, при котором средний силуэт максимален.
\end{itemize}

В данной работе для оценки качества классификации KNN будут использоваться accuracy, precision, recall и F1-мера, а также матрица ошибок. Для выбора оптимального числа кластеров K-Means будет применён метод локтя и оценка силуэта.

\subsection{Краткая информация об инструментах разработки}
% ===== НОВАЯ ПОДГЛАВА (заглушка) =====
\textit{Текст будет написан после получения модели.}

\subsection{Описание источников данных}
% ===== НОВАЯ ПОДГЛАВА (заглушка) =====
\textit{Текст будет написан после получения модели.}



% ========== ГЛАВА 2 ==========
\section{Построение моделей машинного обучения и анализ результатов}

\subsection{Получение данных для построения моделей}
данные Датасет

\subsection{Предобработка и исследовательский анализ данных}
% ===== ЗАГЛУШКА =====
\textit{Текст будет написан после получения модели.}

\subsection{Построение моделей}
сама моделька

\subsection{Сравнительный анализ и оценка качества работы моделей}
% ===== ЗАГЛУШКА =====
\textit{Текст будет написан после получения модели.}  

\newpage

% ========== ЗАКЛЮЧЕНИЕ ==========
\section*{ЗАКЛЮЧЕНИЕ}
\addcontentsline{toc}{section}{ЗАКЛЮЧЕНИЕ}

Здесь будут выводы

\newpage
\begin{thebibliography}{99}

\bibitem{blagov2020}
Благов А.В., Рыцарев И.А. Анализ социальных сетей: учебное пособие. — Самара: Издательство Самарского университета, 2020. — 102 с.

\bibitem{savin2020}
Савин И.В., Мариев О.С., Пушкарев А.А. Методы анализа социальных сетей в экономике: учебное пособие; под общ. ред. И.В. Савина. — Екатеринбург: Издательство Уральского университета, 2020. — 99 с.

\bibitem{vorontsov}
Воронцов К.В. Машинное обучение (курс лекций). — URL: http://www.machinelearning.ru (дата обращения: 18.05.2025).

\bibitem{chubukova2024}
Чубукова И.А. Data Mining: учебное пособие. — 4-е изд. — М.: ИНФРА-М, 2024. — 204 с.

\bibitem{zhdanova2025}
Жданова А.Н. Онтологический подход к предсказанию психологических характеристик пользователей соцсетей на основе графов знаний // Искусственный интеллект и принятие решений. — 2025. — № 3. — С. 40–59.

\bibitem{ratushnyak2022}
Ратушняк Г.Я., Золкин А.Л. Технологии разработки и проектирования информационных систем. Часть 1: учебное пособие. — Москва: Русайнс, 2022. — 201 с.

\bibitem{zaramenskikh2025}
Зараменских Е.П. Разработка информационных систем: учебник и практикум. — 2-е изд. — Москва: Юрайт, 2025. — 78 с.

\end{thebibliography}

\end{document}